\subsection{Non-graph Adaptive Question Selection}

The non-graph learner predicts alternative choices from respondent attributes,
the current survey context, the attributes and revealed answers of previous
questions, and the attributes of each candidate choice task. A shared event MLP
encodes each observed question--answer pair. Mean pooling produces an
order-invariant history representation, which is combined with nine respondent
covariates, the context, and history-length indicators. A shared utility head
scores alternatives A, B, C, and NONE, followed by a softmax. Neither respondent
identifiers nor question identifiers enter the network as features. We average
the predicted probabilities of five respondent-bootstrap models.

\paragraph{Offline training.}
We exclude the union of respondents in all ten evaluation replicates from both
training and validation. The learner uses 1,200 additional respondents for
training and 300 for validation, with all contexts of a respondent kept in the
same partition. Training episodes contain between zero and eight observed
answers; histories of length at least three include the prescribed initial
questions. The remaining history questions are sampled without replacement.
We minimize cross-entropy on answers outside the observed history and select
checkpoints using validation log-loss. Each ensemble member is trained for at
most 120 epochs with 1,024 episodes per epoch and early-stopping patience 30.
MNL and panel MMNL responses are used to train separate ensembles. Network
weights remain fixed throughout the evaluation surveys.

\paragraph{Acquisition rule.}
Let $h$ denote the revealed history and $\mathcal C$ the remaining candidate
questions. For a candidate $q$, we consider its four possible answers $a$ and
their current predicted probabilities $w_a(q)=p_\theta(a\mid q,h)$.
For another candidate $r$, let $p_r^{q,a}$ be its predicted choice distribution
after appending the hypothetical observation $(q,a)$ to $h$. We select the
question maximizing
\begin{equation}
\begin{aligned}
 J(q) &= \frac{1}{|\mathcal C|-1}
     \sum_{r\in\mathcal C\setminus\{q\}}\left[
       H\!\left(\sum_a w_a(q)p_r^{q,a}\right)\right.\\[-2pt]
 &\hspace{35mm}\left.-\sum_a w_a(q)H(p_r^{q,a})\right].
\end{aligned}
\end{equation}
This weighted Jensen--Shannon score measures changes in predictions for other
remaining questions. Its value is used as a predictive-dependence proxy;
coefficient recovery is assessed separately by the common final estimator.
Hypothetical answers never modify the actual respondent history.

\paragraph{Survey collection.}
Each replicate contains 500 respondents in each of the three contexts. All
methods use the supplied context assignments and the same three initial
questions and answers. Non-graph selection then asks six additional questions
from the fifteen remaining questions in the same context. Fixed selection uses
the supplied order of the original nine-question block. Answers are read from
the existing response bank only after a question is selected. The ten
replicates resample respondents from this fixed bank. All nine revealed answers,
including the initial three, are retained for the common MNL or panel MMNL
coefficient-recovery analysis. No prediction holdout is used in this experiment.
